\chapter{Probability at Speed}\label{ch:qm:probability-at-speed}

At ten to four the imbalance of the closing auction appears on the screen, and the desk's
running estimate of the closing price, a number it has been revising all day, jumps by 14
cents. A researcher asks a sharper question than whether the estimate was right: could its
revisions have been predicted from the revisions before them? If they could, the estimate was
not using what it knew, and someone trading against it would have made money. An honest
forecast of a fixed quantity is a \omterm{def:qm:probability-at-speed:condexp}{conditional expectation} along a flow of information, and
such a process has one defining property: its revisions cannot be forecast. This chapter sets
out, at the speed of a reader who has met measure theory, the objects the series computes
with (information, \omterm{def:qm:probability-at-speed:condexp}{conditional expectation}, \omterm{def:qm:probability-at-speed:martingale}{martingales}, \omterm{def:qm:probability-at-speed:stopping}{stopping times}, changes of measure)
and closes by fixing the notation every later book uses.

\section{Information: sigma-algebras, filtrations and conditional expectation}

A probability space $(\Omega, \mathcal F, \P)$ is taken as known. What finance adds is
time: at each instant some events are decided, and the decided ones grow.

\begin{definition}[Filtration, adapted process]\label{def:qm:probability-at-speed:filtration}
A \emph{filtration}\index{filtration} is a family $\mathbb F = (\mathcal F_t)_{t \in
\mathbb T}$ of sub-$\sigma$-algebras of $\mathcal F$, increasing in $t$: $\mathcal F_s
\subseteq \mathcal F_t$ for $s \le t$. The index set $\mathbb T$ is $\{0, 1, \dots, n\}$ or
$[0, T]$; in continuous time the filtration is assumed right-continuous and complete (the
\emph{usual conditions}). A process $(X_t)$ is an \emph{adapted process}\index{adapted
process} if each $X_t$ is $\mathcal F_t$-measurable: its value at $t$ is known at $t$.
\end{definition}

$\mathcal F_t$ is what is known at $t$: the prices printed so far, the cards turned over. A
trading rule must be adapted; a backtest that is not has looked into the future (chapter~3).

\begin{definition}[Conditional expectation]\label{def:qm:probability-at-speed:condexp}
Let $X$ be integrable and $\mathcal G \subseteq \mathcal F$ a sub-$\sigma$-algebra. The
\emph{conditional expectation}\index{conditional expectation} $\E[X \mid \mathcal G]$ is
the almost surely unique $\mathcal G$-measurable, integrable random variable $Y$ with
$\E[X\mathbf 1_A] = \E[Y\mathbf 1_A]$ for every $A \in \mathcal G$. We write $\E_t[X] =
\E[X \mid \mathcal F_t]$.
\end{definition}

Existence is the Radon--Nikodym theorem applied to the measure $A \mapsto \E[X\mathbf
1_A]$ on $\mathcal G$. For square-integrable $X$ there is a more useful picture:
$\E[X \mid \mathcal G]$ is the orthogonal projection of $X$ onto the closed subspace
$L^2(\mathcal G)$, the best forecast of $X$ in mean square among all $\mathcal
G$-measurable random variables. Everything a forecaster can compute from the information
$\mathcal G$ is in that subspace, and the forecast error $X - \E[X \mid \mathcal G]$ is
orthogonal to all of it.

\begin{proposition}[Rules of conditional expectation]\label{prop:qm:probability-at-speed:rules}
For integrable $X, Y$ and $\mathcal H \subseteq \mathcal G$:
(i) \emph{tower}: $\E[\E[X \mid \mathcal G] \mid \mathcal H] = \E[X \mid \mathcal H]$;
(ii) \emph{taking out what is known}: if $Y$ is $\mathcal G$-measurable and $XY$ integrable,
$\E[XY \mid \mathcal G] = Y\E[X \mid \mathcal G]$;
(iii) \emph{independence}: if $X$ is independent of $\mathcal G$, $\E[X \mid \mathcal G] =
\E[X]$;
(iv) \emph{Jensen}: for convex $\varphi$ with $\varphi(X)$ integrable, $\varphi(\E[X \mid
\mathcal G]) \le \E[\varphi(X) \mid \mathcal G]$.
\end{proposition}

\begin{proof}
Check the defining integrals: (i) for $A \in \mathcal H \subseteq \mathcal G$, $\E[\E[X \mid
\mathcal G]\mathbf 1_A] = \E[X\mathbf 1_A]$; (ii) for $Y = \mathbf 1_B$, $B \in \mathcal G$,
then limits; (iii) $\E[X\mathbf 1_A] = \E[X]\P(A)$; (iv) $\varphi$ is a supremum of
countably many affine functions.
\end{proof}

\Cref{fig:qm:probability-at-speed:tree} shows the rules at work on the smallest example
that has them all. A price starts at 100 and moves one dollar up or down on each of two
days with equal probability; $X = (S_2 - 100)^+$. After one day the information is which
half of the tree we are in, and $\E_1[X]$ averages $X$ over that half.

\begin{omfigure}
\begin{tikzpicture}[x=3.2cm,y=0.95cm,font=\small,
  nd/.style={draw=omDef,rounded corners=2pt,fill=omDef!6,inner sep=3pt,align=center}]
\node[nd] (a) at (0,0) {$S_0 = 100$\\$\E_0[X] = 0.5$};
\node[nd] (b) at (1,1.6) {$S_1 = 101$\\$\E_1[X] = 1$};
\node[nd] (c) at (1,-1.6) {$S_1 = 99$\\$\E_1[X] = 0$};
\node[nd] (d) at (2,2.4) {HH: $S_2 = 102$, $X = 2$};
\node[nd] (e) at (2,0.8) {HT: $S_2 = 100$, $X = 0$};
\node[nd] (f) at (2,-0.8) {TH: $S_2 = 100$, $X = 0$};
\node[nd] (g) at (2,-2.4) {TT: $S_2 = 98$, $X = 0$};
\draw[-{Stealth},thick,gray] (a) -- (b); \draw[-{Stealth},thick,gray] (a) -- (c);
\draw[-{Stealth},thick,gray] (b) -- (d.west); \draw[-{Stealth},thick,gray] (b) -- (e.west);
\draw[-{Stealth},thick,gray] (c) -- (f.west); \draw[-{Stealth},thick,gray] (c) -- (g.west);
\draw[omThm,dashed,rounded corners] (2.72,3.0) rectangle (3.02,0.2);
\draw[omThm,dashed,rounded corners] (2.72,-0.2) rectangle (3.02,-3.0);
\node[omThm,font=\footnotesize,align=left,anchor=west] at (3.05,1.6) {an atom\\of $\mathcal F_1$};
\node[omThm,font=\footnotesize,align=left,anchor=west] at (3.05,-1.6) {an atom\\of $\mathcal F_1$};
\node[font=\footnotesize,gray] at (0,-3.4) {$t = 0$};
\node[font=\footnotesize,gray] at (1,-3.4) {$t = 1$};
\node[font=\footnotesize,gray] at (2,-3.4) {$t = 2$};
\end{tikzpicture}

\omcaption{\omterm{def:qm:probability-at-speed:condexp}{Conditional expectation} as averaging over what is still unknown. $\mathcal F_1$
splits the four outcomes into two atoms (dashed); $\E_1[X]$ is the average of $X$ on the
atom reached, and $\E_0[X]$ is the average of $\E_1[X]$: the tower rule. The three values
$0.5 \to 1 \to 2$ along the top path form a \omterm{def:qm:probability-at-speed:doob}{Doob martingale}.}\label{fig:qm:probability-at-speed:tree}
\end{omfigure}

Laws are handled through their transforms. The one the series uses most is defined here, so
that later books (Lévy processes in chapter~6, Fourier pricing in chapter~28, stochastic
volatility in One Quant Book~5) can point to it.

\begin{definition}[Characteristic function]\label{def:qm:probability-at-speed:charfn}
The \emph{characteristic function}\index{characteristic function} of a random variable $X$
is $\varphi_X(u) = \E[e^{\iu uX}]$, $u \in \R$; of a random vector, $\varphi_X(u) =
\E[e^{\iu u^\top X}]$, $u \in \R^d$. It exists for every law, determines it, and turns sums
of independent variables into products.
\end{definition}

\begin{theorem}[Central limit theorem]\label{thm:qm:probability-at-speed:clt}
Let $X_1, X_2, \dots$ be independent with means $\mu_i$ and variances $\sigma_i^2$, $s_n^2 =
\sum_{i \le n}\sigma_i^2$, and suppose Lindeberg's condition $s_n^{-2}\sum_{i\le n}\E[(X_i -
\mu_i)^2\mathbf 1_{\{|X_i - \mu_i| > \varepsilon s_n\}}] \to 0$ holds for every $\varepsilon >
0$ (it does for identically distributed variables with finite variance). Then
$s_n^{-1}\sum_{i \le n}(X_i - \mu_i) \xrightarrow{d} \mathcal N(0, 1)$.
\end{theorem}

\admitted

The proof expands the \omterm{def:qm:probability-at-speed:charfn}{characteristic function} to second order and uses Lévy's continuity
theorem (Williams, 1991). Every standard error in Part~II rests on it.

\section{Martingales and the Doob martingale of a forecast}

\begin{definition}[Martingale]\label{def:qm:probability-at-speed:martingale}
An adapted, integrable process $(M_t)$ is a \emph{martingale}\index{martingale} if
$\E_s[M_t] = M_s$ for all $s \le t$; a \emph{submartingale}\index{submartingale} if
$\E_s[M_t] \ge M_s$; a \emph{supermartingale}\index{supermartingale} if $\E_s[M_t] \le
M_s$.
\end{definition}

A \omterm{def:qm:probability-at-speed:martingale}{martingale} is a fair game seen from any date. Four examples recur in the book. With $\xi_i$ independent, taking the values $\pm 1$ with
probabilities $p$ and $q = 1 - p$, and $S_n = \sum_{i \le n}\xi_i$:
\begin{itemize}
\item $S_n$ is a \omterm{def:qm:probability-at-speed:martingale}{martingale} when $p = \tfrac12$, a \omterm{def:qm:probability-at-speed:martingale}{submartingale} when $p > \tfrac12$;
\item $S_n^2 - n$ is a \omterm{def:qm:probability-at-speed:martingale}{martingale} when $p = \tfrac12$: the variance grows by one per step;
\item $(q/p)^{S_n}$ is a \omterm{def:qm:probability-at-speed:martingale}{martingale} for every $p$, since $\E[(q/p)^{\xi}] = p\,q/p + q\,p/q = 1$;
\item $\exp(\theta S_n - n\ln\cosh\theta)$ is a \omterm{def:qm:probability-at-speed:martingale}{martingale} when $p = \tfrac12$, for every real
$\theta$: the exponential \omterm{def:qm:probability-at-speed:martingale}{martingale}, prototype of the density processes of
\cref{sec:qm:probability-at-speed:measure}.
\end{itemize}

The example that matters most for a trading desk is not a price but a forecast.

\begin{definition}[Doob martingale]\label{def:qm:probability-at-speed:doob}
For an integrable random variable $X$ and a \omterm{def:qm:probability-at-speed:filtration}{filtration} $\mathbb F$, the \emph{Doob
martingale}\index{Doob martingale} of $X$ is $M_t = \E_t[X]$.
\end{definition}

It is a \omterm{def:qm:probability-at-speed:martingale}{martingale} by the tower rule. Every honest forecast of a fixed quantity (a closing
price, the sum of five cards, next month's inflation print) is a \omterm{def:qm:probability-at-speed:doob}{Doob martingale} along the
forecaster's \omterm{def:qm:probability-at-speed:filtration}{filtration}, and what makes it testable is that its revisions are orthogonal.

\begin{proposition}[Orthogonal increments]\label{prop:qm:probability-at-speed:orthogonal}
Let $(M_k)_{k = 0,\dots,n}$ be a square-integrable \omterm{def:qm:probability-at-speed:martingale}{martingale} with increments $D_k = M_k -
M_{k-1}$. Then $\E[D_k Y] = 0$ for every square-integrable $\mathcal F_{k-1}$-measurable
$Y$; in particular $\E[D_jD_k] = 0$ for $j < k$, and
\[
\Var(M_n) = \Var(M_0) + \sum_{k=1}^n \E[D_k^2].
\]
\end{proposition}

\begin{proof}
$\E[D_kY] = \E[Y\E_{k-1}[D_k]] = 0$ by taking out what is known. With $Y = D_j$, $j < k$, the
cross terms of $(M_n - M_0)^2 = (\sum_k D_k)^2$ vanish, and $\Cov(M_0, D_k) = 0$ likewise.
\end{proof}

Two consequences drive this chapter's tutorial and problem. First, a forecast whose
revisions can be predicted from the past (a regression of $D_k$ on $D_{k-1}$ with a nonzero
slope) is not a \omterm{def:qm:probability-at-speed:condexp}{conditional expectation}: it is leaving information unused. Second, the
variance of the final value splits into the variances of the revisions, so one can say how
much of a closing price's uncertainty is resolved at each time of day
(\cref{fig:qm:probability-at-speed:day}).

\begin{omfigure}
\begin{tikzpicture}
\begin{groupplot}[group style={group size=2 by 1,horizontal sep=1.6cm},
  width=7.6cm,height=4.8cm,xlabel={minutes after 09:30},
  tick label style={font=\small},label style={font=\small},
  xmin=0,xmax=395,grid=major,grid style={gray!25},xtick={0,120,240,380},
  table/col sep=comma]
\nextgroupplot[ylabel={forecast of the close (USD)},ymin=99.2,ymax=102.4]
\addplot[omDef,thick] table[x=minute,y=forecast]{figdata/methods/01-probability-at-speed/day.csv};
\draw[omThm,dashed] (axis cs:380,99.2) -- (axis cs:380,102.4);
\node[font=\footnotesize,anchor=east,omThm] at (axis cs:374,99.5) {15:50 news: $+0.48$};
\nextgroupplot[ylabel={share of Var(close) resolved},ymin=0,ymax=1.05,
  legend style={font=\footnotesize,at={(0.03,0.97)},anchor=north west,fill=white,draw=none}]
\addplot[omDef,very thick] table[x=minute,y=theory]{figdata/methods/01-probability-at-speed/resolved.csv};
\addplot[omThm,only marks,mark=*,mark size=1.3pt] table[x=minute,y=sim]{figdata/methods/01-probability-at-speed/resolved.csv};
\legend{formula,20\,000 days}
\end{groupplot}
\end{tikzpicture}

\omcaption{Left: one simulated day of $M_t = \E_t[\text{close}]$ in the weekend problem's
model (a random walk with a daily standard deviation of USD~1.80 on a USD~100 stock, and news
on the closing imbalance with standard deviation USD~0.30 at 15:50, $+0.48$ on this day).
Right: the share of the close's variance resolved by each time of day, by
\cref{prop:qm:probability-at-speed:orthogonal}: linear through the session, a step of 2.7\%
at the news, the last 2.5\% in the final ten minutes. Data: the chapter's tutorial,
seeded.}\label{fig:qm:probability-at-speed:day}
\end{omfigure}

The maximal inequality bounds how far a \omterm{def:qm:probability-at-speed:martingale}{martingale} wanders over a period by where it ends;
chapter~2 uses it to control Brownian paths.

\begin{proposition}[Doob's maximal inequality]\label{prop:qm:probability-at-speed:doob}
For a nonnegative \omterm{def:qm:probability-at-speed:martingale}{submartingale} $(X_k)_{k \le n}$ and $c > 0$, $c\,\P(\max_{k\le n}X_k \ge c)
\le \E[X_n\mathbf 1_{\{\max_k X_k \ge c\}}] \le \E[X_n]$; and for $p > 1$, $\E[\max_{k \le
n}X_k^p] \le (p/(p-1))^p\E[X_n^p]$. For a square-integrable \omterm{def:qm:probability-at-speed:martingale}{martingale}, applied to $|M_k|$:
$\E[\max_{k\le n}M_k^2] \le 4\E[M_n^2]$.
\end{proposition}

\admitted

\section{Stopping times and optional stopping}

\begin{definition}[Stopping time]\label{def:qm:probability-at-speed:stopping}
A random time $\tau$ with values in $\mathbb T \cup \{\infty\}$ is a \emph{stopping
time}\index{stopping time} if $\{\tau \le t\} \in \mathcal F_t$ for every $t$: whether it
has occurred by $t$ is known at $t$.
\end{definition}

``The first time the price touches 101'' is a \omterm{def:qm:probability-at-speed:stopping}{stopping time}; ``the time of the day's high''
is not, since it is known only at the close. When is a \omterm{def:qm:probability-at-speed:martingale}{martingale}'s expected value at a
\omterm{def:qm:probability-at-speed:stopping}{stopping time} still its starting value?

\begin{definition}[Uniform integrability]\label{def:qm:probability-at-speed:ui}
A family $(X_i)_{i \in I}$ of random variables has \emph{uniform
integrability}\index{uniform integrability} if $\sup_{i}\E[|X_i|\mathbf 1_{\{|X_i| > c\}}]
\to 0$ as $c \to \infty$.
\end{definition}

\begin{theorem}[Optional stopping]\label{thm:qm:probability-at-speed:optional}
Let $(M_n)$ be a \omterm{def:qm:probability-at-speed:martingale}{martingale} and $\tau$ a \omterm{def:qm:probability-at-speed:stopping}{stopping time}. Then $\E[M_\tau] = \E[M_0]$ if either
(i) $\tau$ is bounded, or (ii) $\tau < \infty$ almost surely and the stopped process
$(M_{n \wedge \tau})$ is uniformly integrable; in particular if $|M_{n\wedge\tau}| \le c$ for
all $n$.
\end{theorem}

\begin{proof}
(i) If $\tau \le N$, then $M_\tau = M_0 + \sum_{k=1}^N D_k\mathbf 1_{\{\tau \ge k\}}$, and
$\{\tau \ge k\} = \{\tau \le k-1\}^c \in \mathcal F_{k-1}$, so each term has zero mean by
\cref{prop:qm:probability-at-speed:orthogonal}. (ii) Apply (i) to $\tau \wedge n$; $M_{n \wedge
\tau} \to M_\tau$ almost surely, and \omterm{def:qm:probability-at-speed:ui}{uniform integrability} upgrades this to convergence in
$L^1$ (Vitali), so the means converge.
\end{proof}

The theorem prices every take-profit and stop-loss rule on a \omterm{def:qm:probability-at-speed:martingale}{martingale}. A position whose
value moves by $\pm 1$ tick with equal probability is closed at $+a$ or $-b$. The stopped
value is bounded, so $\E[S_\tau] = 0$: with $\P(\text{up first}) = \pi$,
$\pi a - (1 - \pi)b = 0$ and
\[
\pi = \frac{b}{a+b}.
\]
A stop close to entry and a distant profit target win rarely and lose often, and the
expected profit is zero whatever $a$ and $b$ are: the rule reshapes the distribution of the
result, not its mean. With a drift ($p \ne \tfrac12$) the \omterm{def:qm:probability-at-speed:martingale}{martingale} $(q/p)^{S_n}$ gives
$\pi = (1 - (q/p)^b)/(1 - (q/p)^{a+b})$; \cref{fig:qm:probability-at-speed:stops} compares
both formulas with simulation.

\begin{omfigure}
\begin{tikzpicture}
\begin{axis}[width=11cm,height=5cm,
  xlabel={stop-loss distance $b$ (ticks), take-profit $a = 10$},ylabel={P(take-profit first)},
  tick label style={font=\small},label style={font=\small},
  xmin=0,xmax=21,ymin=0,ymax=0.75,grid=major,grid style={gray!25},
  xtick={2,4,6,8,10,12,14,16,18,20},
  legend columns=2,legend style={font=\footnotesize,at={(0.5,-0.3)},anchor=north,draw=none,fill=none,/tikz/every even column/.append style={column sep=8pt}},
  table/col sep=comma]
\addplot[omDef,very thick] table[x=b,y=fair_theory]{figdata/methods/01-probability-at-speed/stops.csv};
\addplot[omDef,only marks,mark=*,mark size=1.8pt] table[x=b,y=fair_sim]{figdata/methods/01-probability-at-speed/stops.csv};
\addplot[omThm,very thick] table[x=b,y=drift_theory]{figdata/methods/01-probability-at-speed/stops.csv};
\addplot[omThm,only marks,mark=triangle*,mark size=2.2pt] table[x=b,y=drift_sim]{figdata/methods/01-probability-at-speed/stops.csv};
\legend{{$p = 0.5$: $b/(a+b)$},simulated,{$p = 0.49$: formula},simulated}
\end{axis}
\end{tikzpicture}

\omcaption{Probability that a $\pm1$-tick random walk reaches a take-profit 10 ticks away
before a stop-loss $b$ ticks away: the optional-stopping formulas (lines) and 20\,000 seeded
paths per point (marks). A 1-point drift against the position ($p = 0.49$) costs 10 points
of probability at $b = 10$. Data: the chapter's tutorial.}\label{fig:qm:probability-at-speed:stops}
\end{omfigure}

\begin{remark}[The doubling strategy]\label{rem:qm:probability-at-speed:doubling}
Doubling the stake after each loss of a fair even-money bet and stopping at the first win
yields $+1$ almost surely, apparently contradicting the theorem. The \omterm{def:qm:probability-at-speed:stopping}{stopping time} is finite
but unbounded, and the stopped process is not uniformly integrable: before the win it has
lost $2^k - 1$ with probability $2^{-k}$. A finite credit line bounds the stakes, makes
condition (ii) hold, and restores $\E[M_\tau] = 0$: a small, frequent gain paid for by a rare,
enormous loss.
\end{remark}

\section{Changing the measure}\label{sec:qm:probability-at-speed:measure}

Much of quantitative finance computes under a measure other than the one that describes the
world: one under which discounted prices are \omterm{def:qm:probability-at-speed:martingale}{martingales} (chapter~5), or one under which a
rare loss is common enough to simulate (chapter~26).

\begin{definition}[Change of measure]\label{def:qm:probability-at-speed:change}
Two probability measures $\P$ and $\mathbb Q$ on $(\Omega, \mathcal F)$ are
\emph{equivalent measures}\index{equivalent measures} if they have the same null sets. A
\emph{change of measure}\index{change of measure} from $\P$ to an equivalent $\mathbb Q$ is
described by the \emph{Radon--Nikodym derivative}\index{Radon--Nikodym derivative} $Z =
d\mathbb Q/d\P$, the almost surely unique positive random variable with $\mathbb Q(A) =
\E[Z\mathbf 1_A]$ for all $A \in \mathcal F$; then $\E^{\mathbb Q}[X] = \E[ZX]$. Along a
\omterm{def:qm:probability-at-speed:filtration}{filtration}, the \emph{density process}\index{density process} is $Z_t = \E_t[Z]$, the
Radon--Nikodym derivative of $\mathbb Q$ restricted to $\mathcal F_t$.
\end{definition}

The \omterm{def:qm:probability-at-speed:change}{density process} is a positive $\P$-martingale with $Z_0 = 1$ (a \omterm{def:qm:probability-at-speed:doob}{Doob martingale}), and
conversely every such \omterm{def:qm:probability-at-speed:martingale}{martingale} defines a measure. \omterm{def:qm:probability-at-speed:condexp}{Conditional expectations} under the new
measure follow from the old ones.

\begin{proposition}[Bayes formula for a change of measure]\label{prop:qm:probability-at-speed:bayes}
If $\mathbb Q \sim \P$ with \omterm{def:qm:probability-at-speed:change}{density process} $(Z_t)$ and $X$ is $\mathbb Q$-integrable, then
for $s \le t$ and $X$ $\mathcal F_t$-measurable,
\[
\E^{\mathbb Q}_s[X] = \frac{\E_s[Z_tX]}{Z_s}.
\]
In particular, an adapted $(X_t)$ is a $\mathbb Q$-martingale if and only if $(Z_tX_t)$ is a
$\P$-martingale.
\end{proposition}

\begin{proof}
The right-hand side $Y$ is $\mathcal F_s$-measurable. For $A \in \mathcal F_s$,
\[
\E^{\mathbb Q}[Y\mathbf 1_A] = \E[Z_sY\mathbf 1_A] = \E\bigl[\E_s[Z_tX]\mathbf 1_A\bigr] =
\E[Z_tX\mathbf 1_A] = \E^{\mathbb Q}[X\mathbf 1_A],
\]
using $\E[Z W] = \E[Z_s W]$ for $\mathcal F_s$-measurable $W$, and the tower rule.
\end{proof}

\begin{example}[Shifting a Gaussian, and seeing a far tail]\label{ex:qm:probability-at-speed:shift}
Let $X \sim \mathcal N(0, 1)$ under $\P$ and $Z = \exp(cX - c^2/2)$. Then $\E^{\mathbb
Q}[e^{\iu uX}] = \E[e^{(c + \iu u)X - c^2/2}] = e^{\iu uc - u^2/2}$: under $\mathbb Q$, $X \sim
\mathcal N(c, 1)$. The \omterm{def:qm:probability-at-speed:change}{change of measure} has moved the mean without touching the shape. Run it
backwards to estimate $p = \P(X > 4) = 3.17 \times 10^{-5}$: sample $Y \sim \mathcal N(4, 1)$
and average $e^{-4Y + 8}\mathbf 1_{\{Y > 4\}}$, the indicator reweighted by $d\P/d\mathbb Q$.
With 100\,000 draws, plain sampling sees three exceedances and has a standard error of $1.7
\times 10^{-5}$, half the answer; the reweighted estimate has a standard error of $2.1 \times
10^{-7}$, eighty times smaller. Chapter~26 turns this into importance sampling, and chapter~5
does the same computation for whole Brownian paths.
\end{example}

\section{The notation of the series}\label{sec:qm:probability-at-speed:notation}

Every book of the series from this one on uses the symbols below. They extend the tables of
One Quant Books 1 and 2, whose meanings are kept: $s_t$ is the bid--ask spread, $\mathcal S$
a credit spread, $\tau = T - t$ a time to expiry, $r$ a rate, $\ell$ a borrow fee, $K$ a
strike, $P(t,T)$ a discount factor, $y$ a yield, $\delta$ an accrual fraction and $R$ a
recovery rate. A symbol a chapter declares \emph{local} may carry another meaning there.

\begin{notation}[Probability, measures and processes]\label{not:qm:probability-at-speed:prob}
\leavevmode\par\smallskip\noindent
{\small
\begin{tabular}{@{}p{4.3cm}p{9.7cm}@{}}
\toprule
$(\Omega,\mathcal F,\P)$, $\mathbb F = (\mathcal F_t)$ & probability space; \omterm{def:qm:probability-at-speed:filtration}{filtration} (usual conditions) \\
$\E_t[X] = \E[X \mid \mathcal F_t]$, $\E^{\mathbb Q}_t$ & \omterm{def:qm:probability-at-speed:condexp}{conditional expectation}; under another measure \\
$\P$, $\mathbb Q$, $B_t = e^{\int_0^t r_s\,ds}$ & real-world measure; risk-neutral measure; money-market account \\
$\mathcal N_t$, $\mathbb Q^{\mathcal N}$; $\mathbb Q^T$, $\mathbb Q^A$ & generic numeraire and its measure; $T$-forward and annuity measures \\
$d\mathbb Q/d\P$, $Z_t$ & \omterm{def:qm:probability-at-speed:change}{Radon--Nikodym derivative}; \omterm{def:qm:probability-at-speed:change}{density process} \\
$\Phi$, $\varphi$; $\mathcal N(m, s^2)$ & standard normal cdf and pdf (never $N(d_1)$); the normal law, always with arguments \\
$\varphi_X(u) = \E[e^{\iu uX}]$ & \omterm{def:qm:probability-at-speed:charfn}{characteristic function}, always subscripted \\
$\mathbf 1_A$; $\overset{d}{=}$, $\xrightarrow{d}$, $\xrightarrow{\P}$ & indicator; equality and convergence in law, in probability \\
$\tau$ & a \omterm{def:qm:probability-at-speed:stopping}{stopping time}; where it meets a time to expiry, write $T - t$; default times are subscripted ($\tau_C$) \\
$W_t$; $W^{\mathbb Q}_t$ & Brownian motion under the measure in force; decorated when two measures appear \\
$[X]_t$, $[X,Y]_t$, $\mathcal E(X)_t$ & quadratic variation, covariation, stochastic exponential \\
$dX = \mu\,dt + \sigma\,dW$; $\mathcal L$, $\mathcal L^*$ & a diffusion; its generator and the adjoint \\
$dX = \kappa(\bar x - X)\,dt + \sigma\,dW$ & Ornstein--Uhlenbeck: speed $\kappa$, level $\bar x$, half-life $\ln 2/\kappa$ \\
$dv = \kappa(\bar v - v)\,dt + \eta\sqrt v\,dW$ & square-root process; Feller condition $2\kappa\bar v \ge \eta^2$; $\eta$ is the vol-of-vol throughout \\
$N_t$, $t_1 < t_2 < \dots$ & counting process and its event times \\
$\lambda_t$, $\Lambda_t = \int_0^t\lambda_s\,ds$ & intensity (Poisson rate, hazard rate, Hawkes intensity $\lambda_t = \mu + \sum_{t_i < t}g(t - t_i)$); compensator \\
$(\sigma^2, \nu, \gamma)$, $\psi$ & Lévy triplet and characteristic exponent, $\varphi_{X_t}(u) = e^{t\psi(u)}$ \\
$H$, $W^H_t$ & Hurst exponent; fractional Brownian motion \\
\bottomrule
\end{tabular}}
\end{notation}

\begin{notation}[Statistics, matrices and numerics]\label{not:qm:probability-at-speed:stat}
\leavevmode\par\smallskip\noindent
{\small
\begin{tabular}{@{}p{4.3cm}p{9.7cm}@{}}
\toprule
$n$, $\theta \in \Theta$, $\theta_0$ & sample size, parameter, true value; hat an estimate, tilde a shrunk estimate, bar a sample mean \\
$\ell_n(\theta)$, $\mathcal I(\theta)$ & log-likelihood (always with subscript and argument); Fisher information \\
$\mathrm{se}(\hat\theta)$; $\mathrm{SR}$, $\widehat{\mathrm{SR}}$ & standard error; Sharpe ratio and its estimate \\
$R_t = \ln(S_t/S_{t-1})$ & log return ($R$ alone stays the recovery rate) \\
$L$, $\Delta = 1 - L$; $\gamma(h)$, $\rho(h)$ & lag and difference operators; autocovariance, autocorrelation \\
$\phi_i$, $\vartheta_j$, $\varepsilon_t$ & autoregressive and moving-average coefficients (always indexed); innovations \\
$\mathrm{RV}_t$, $\mathrm{IV}_t$; $\sigma_{\mathrm{imp}}$ & realised and integrated variance; implied volatility (never ``IV'' in a formula) \\
$\Sigma$, $C$, $\hat\Sigma$; $\mathbf 1$, $I_n$, ${}^\top$ & covariance, correlation, sample covariance; ones vector, identity, transpose; vectors are columns \\
$\lambda_1 \ge \dots \ge \lambda_N$, $\operatorname{cond}(A)$ & eigenvalues (always indexed); condition number \\
$\min f(x)$ s.t.\ $g_i(x) \le 0$, $h_j(x) = 0$ & optimisation; multipliers $u \ge 0$, $v$; optimum $x^\star$, values $p^\star$, $d^\star$ \\
$\mathrm{fl}(x)$, $u = 2^{-53}$, $\varepsilon_{\mathrm{mach}} = 2^{-52}$ & floating point in binary64, round to nearest; $\mathrm{ulp}(x)$ \\
$t_k = k\Delta t$, $x_j$, $V^k_j$; $M$, $\hat V_M$ & time grid, space grid, numerical solution; Monte Carlo paths and estimator \\
\bottomrule
\end{tabular}}
\end{notation}

The options, rates, credit and risk symbols ($\sigma_{\mathrm{imp}}(K,T)$, $k = \ln(K/F_{0,T})$,
$\xi_t(u)$, the Greeks with the rate Greek written Rho because $\rho$ is always a correlation,
$S_{a,b}(t)$, $A_{a,b}(t)$, $Q_C(t)$, $\mathrm{VaR}_{\alpha,h}$, $\mathrm{ES}_{\alpha,h}$) come
from the same table and are introduced in One Quant Books 5 and 6.

\section{Tutorial: the Doob martingale of a card game}\label{tut:qm:probability-at-speed:cards}

\begin{tutorial}
\textbf{Goal.} Build the \omterm{def:qm:probability-at-speed:doob}{Doob martingale} of the five-card game of One Quant Book~2,
chapter~30 (the sum of five cards dealt from a 52-card deck, aces 1 to kings 13), check that
its revisions behave as \cref{prop:qm:probability-at-speed:orthogonal} says, and catch a
forecaster that under-reacts. \textbf{End state:} the table below and
\cref{fig:qm:probability-at-speed:day,fig:qm:probability-at-speed:stops}.
\begin{tutsteps}
\item \textbf{The forecast.} After $k$ cards with sum $s_k$, the remaining $52 - k$ cards
have sum $364 - s_k$, so $M_k = s_k + (5 - k)(364 - s_k)/(52 - k)$, and $M_0 = 35$. The
under-reacting forecaster moves only a fraction $a$ of the way from 35 until the last card.
\omcode{code/methods/01-probability-at-speed/python/qm_martingales.py}{17}{39}{The Doob martingale of the card game and an under-reacting forecaster.}\label{lst:qm:probability-at-speed:cards}
\item \textbf{The test.} Regress the last revision on the one before it, pooled over deals;
the running project's function does it for any panel of forecast paths.
\omcode{code/firm/mgtest/firm_mgtest.py}{29}{47}{Regression of a revision on an earlier one: zero slope for a martingale.}\label{lst:qm:probability-at-speed:regression}
\item \textbf{Run} \texttt{card\_table()} over 20\,000 seeded deals, then
\texttt{fig\_martingales.py} for the figures.
\end{tutsteps}
\begin{center}
\small
\begin{tabular}{@{}lcc@{}}
\toprule
 & honest forecast & under-reacting ($a = 0.6$) \\
\midrule
mean revision & $0.001$ & $0.001$ \\
slope of last revision on the fourth & $-0.005$ & $0.661$ \ (theory $0.667$) \\
$t$-statistic of the slope & --- & $45.9$ \\
\bottomrule
\end{tabular}
\end{center}
Both forecasts have revisions with mean zero; only the regression tells them apart.
\textbf{What to change next.} Let the forecaster over-react ($a > 1$) and predict the sign of
the slope; replace the card game by the closing-price model of the weekend problem and find
how many days of data the test needs.
\end{tutorial}


\section{Build: martingale diagnostics}\label{bld:qm:probability-at-speed:mgtest}

\begin{build}
\bfield{Purpose} Every forecast the miniature firm publishes (a fair value, an expected
closing price, a predicted fill rate) is checked for the \omterm{def:qm:probability-at-speed:martingale}{martingale} property before anyone
trades on it or against it.
\bfield{Interface} \texttt{revisions(paths)}; \texttt{revision\_regression(paths, lag=1,
col=None)} returning slope, standard error, $t$ and $n$; \texttt{variance\_ratio(increments,
q)}; \texttt{variance\_shares(paths)}; \texttt{martingale\_report(paths)}. A panel has one
row per day or deal and one column per revision time.
\bfield{Rules} No intercept in the revision regression (the mean revision is reported
separately); variance shares are those of the revisions, which for a \omterm{def:qm:probability-at-speed:martingale}{martingale} sum to the
variance of the final value; everything seeded and vectorised.
\bfield{Acceptance tests} \texttt{code/firm/mgtest/tests/}: random walks pass; a forecast
that adds half of the news late is caught with the right slope; shares are uniform for iid
revisions; the variance ratio is one for iid increments and below one for mean-reverting
ones.
\bfield{Stretch} Heteroskedasticity- and autocorrelation-robust standard errors (chapter~11);
a test of the revisions against any $\mathcal F_{k-1}$-measurable signal, not only past
revisions; multiple-testing control across many forecasts (chapter~12).
\end{build}

\begin{omsources}
\item A.~N.~Kolmogorov, \emph{Grundbegriffe der Wahrscheinlichkeitsrechnung}, Springer,
1933: the measure-theoretic axioms.
\item J.~Ville, \emph{Étude critique de la notion de collectif}, Gauthier-Villars, 1939,
where the word martingale enters probability.
\item J.~L.~Doob, \emph{Stochastic Processes}, Wiley, 1953.
\item O.~Nikodym, ``Sur une généralisation des intégrales de M.~J.~Radon'', \emph{Fundamenta
Mathematicae} 15, 1930.
\item D.~Williams, \emph{Probability with Martingales}, Cambridge University Press, 1991: the
proofs admitted here.
\end{omsources}

\section{Exercises}

\begin{exercise}[$\star$]\label{exo:qm:probability-at-speed:1}
In the five-card game, what is the forecast of the sum after the first card turns out to be
a king?
\end{exercise}

\begin{exercise}[$\star$]\label{exo:qm:probability-at-speed:2}
Which of these are \omterm{def:qm:probability-at-speed:stopping}{stopping times} for the \omterm{def:qm:probability-at-speed:filtration}{filtration} of observed prices: (a) the first time
the price is 1\% above the open; (b) the time of the day's low; (c) 15:50; (d) the first
time after 15:50 that the forecast of the close moves by more than 10 cents?
\end{exercise}

\begin{exercise}[$\star$]\label{exo:qm:probability-at-speed:3}
A position on a fair $\pm1$-tick walk is closed at $+1$ or $-3$. What is the probability of
the take-profit, and why is the expected profit still zero?
\end{exercise}

\begin{exercise}[$\star\star$]\label{exo:qm:probability-at-speed:4}
A forecast of the close is revised by independent news of equal variance in each of the 390
minutes of a session, and nothing else. What share of the close's variance is resolved in
the last half hour?
\end{exercise}

\begin{exercise}[$\star\star$]\label{exo:qm:probability-at-speed:5}
Under $\P$, $X \sim \mathcal N(0,1)$, and $d\mathbb Q/d\P = \exp(cX - c^2/2)$. Compute $\E^{\mathbb
Q}[X]$ and $\mathbb Q(X > c)$, and check that $\mathbb Q$ is a probability measure.
\end{exercise}

\begin{exercise}[$\star\star$]\label{exo:qm:probability-at-speed:6}
The under-reacting card forecaster publishes $F_k = 35 + a(M_k - 35)$ for $k \le 4$ and $F_5 =
M_5$. Compute $\E[F_5 - F_4 \mid \mathcal F_4]$ and the slope of the regression of $F_5 - F_4$ on
$F_4 - 35$.
\end{exercise}

\begin{exercise}[$\star\star\star$]\label{exo:qm:probability-at-speed:7}
\emph{Coding.} With \texttt{revision\_regression}, estimate the slope of the last revision
of the under-reacting card forecaster ($a = 0.6$) on the previous revision over 20\,000
seeded deals, and compare it with the answer to \cref{exo:qm:probability-at-speed:6}.
\end{exercise}

\begin{exercise}[$\star\star\star$]\label{exo:qm:probability-at-speed:8}
\emph{Find the flaw.} ``Over a year our fair-value model's revisions averaged zero to four
decimals, so the model is a \omterm{def:qm:probability-at-speed:martingale}{martingale} and nobody can trade against it.'' Correct it.
\end{exercise}

\section{Problem: The Fair Value of the Close}

\begin{problem}[{Weekend problem --- how much of the close is still unknown at ten to four}]\label{pb:qm:probability-at-speed:1}
A desk publishes $M_t = \E_t[C]$, its forecast of a stock's closing price $C$. In its model
the price starts at USD~100 and moves in each of the 390 minutes of the session by
independent Gaussian increments with a total standard deviation of USD~1.80 over the day; at
15:50 the imbalance of the closing auction is published, which moves the forecast by an
independent Gaussian amount $J$ with standard deviation USD~0.30 and zero mean; the last ten
minutes of the session follow.

\medskip\noindent\textbf{Part I --- The forecast.}
\begin{enumerate}
\item What are the variance and standard deviation of $C$?
\item Write $M_t$ before 15:50 and after the news, in terms of the observed price.
\item Why is $(M_t)$ a \omterm{def:qm:probability-at-speed:martingale}{martingale} for the desk's \omterm{def:qm:probability-at-speed:filtration}{filtration}?
\item What is the revision of $M$ at the 15:50 news?
\item What is the probability that the news moves the forecast by more than 14 cents?
\end{enumerate}

\medskip\noindent\textbf{Part II --- What is resolved when.}
\begin{enumerate}[resume]
\item Why do the variances of the revisions add up to $\Var(C)$?
\item What share of $\Var(C)$ is resolved by noon?
\item What share is resolved at the news itself?
\item What share is resolved from the news to the close, inclusive?
\item Check this share against the simulation of 20\,000 days.
\end{enumerate}

\medskip\noindent\textbf{Part III --- A forecast that is too slow.}
\begin{enumerate}[resume]
\item A second desk adds only half of $J$ at 15:50 and the other half at the close. Is its
forecast a \omterm{def:qm:probability-at-speed:martingale}{martingale}?
\item What is the correlation between its 15:50 revision and its revision over the last ten
minutes?
\item What slope does the regression of the second revision on the first give, in theory and
in the simulation?
\item How many days of data does the test need for the slope's $t$-statistic to reach 2?
\item A trader buys one share from that desk at its forecast at 15:50 when its revision is
positive, sells one when it is negative, and closes at the close. What does it earn on average
per share?
\end{enumerate}

\medskip\noindent\textbf{Part IV --- Judgement.}
\begin{enumerate}[resume]
\item If the last ten minutes carry three times the per-minute variance of the rest of the
day (with the day's total unchanged), what share is resolved from the news to the close?
\item Does the answer depend on the price level?
\item Why does the share matter to a trader who can only act in the closing auction?
\item State the \emph{named result}: the share of the close's variance resolved from 15:50.
\item In one sentence: what does the \omterm{def:qm:probability-at-speed:martingale}{martingale} property of a forecast promise, and what
does it not?
\end{enumerate}
\end{problem}

\section{Interview questions}

\begin{interviewq}[$\star$ \iqroles{trader}]\label{iq:qm:probability-at-speed:1}
You play a fair coin game for a dollar a toss and stop as soon as you are one dollar up.
Isn't that a guaranteed win?
\end{interviewq}

\begin{interviewq}[$\star$ \iqroles{researcher}]\label{iq:qm:probability-at-speed:2}
What is a \omterm{def:qm:probability-at-speed:martingale}{martingale}? Is a stock price one?
\end{interviewq}

\begin{interviewq}[$\star\star$ \iqroles{researcher, mle}]\label{iq:qm:probability-at-speed:3}
Show that the increments of a \omterm{def:qm:probability-at-speed:martingale}{martingale} are uncorrelated. What does that imply for a
forecast you publish?
\end{interviewq}

\begin{interviewq}[$\star\star$ \iqroles{trader}]\label{iq:qm:probability-at-speed:4}
At the imbalance publication your forecast of the close moves 14 cents; a competitor's moves
7 cents, then another 7 at the close, every day. What do you do?
\end{interviewq}

\begin{interviewq}[$\star\star$ \iqroles{researcher, risk}]\label{iq:qm:probability-at-speed:5}
When is $\E[X \mid Y]$ the linear regression of $X$ on $Y$?
\end{interviewq}

\begin{interviewq}[$\star\star\star$ \iqroles{developer, researcher}]\label{iq:qm:probability-at-speed:6}
Design a daily check that a real-time fair-value service behaves like a \omterm{def:qm:probability-at-speed:martingale}{martingale}. What can
make the check lie?
\end{interviewq}
